\documentclass[10pt,a4paper]{amsart}
\usepackage{amsmath,amssymb,amsthm,graphicx}
\usepackage{verbatim}
\newenvironment{reference}{\comment}{\endcomment}
%\newenvironment{reference}{}{}
\newenvironment{slogan}{\comment}{\endcomment}
\newenvironment{history}{\comment}{\endcomment}
% For commutative diagrams we use Xy-pic
\usepackage[all]{xy}

% We use 2cell for 2-commutative diagrams.
\xyoption{2cell}
% Theorem environments.
%
\theoremstyle{plain}
\newtheorem{theorem}[subsection]{Theorem}
\newtheorem{proposition}[subsection]{Proposition}
\newtheorem{lemma}[subsection]{Lemma}

\theoremstyle{definition}
\newtheorem{definition}[subsection]{Definition}
\newtheorem{example}[subsection]{Example}
\newtheorem{exercise}[subsection]{Exercise}
\newtheorem{situation}[subsection]{Situation}

\theoremstyle{remark}
\newtheorem{remark}[subsection]{Remark}
\newtheorem{remarks}[subsection]{Remarks}

\numberwithin{equation}{subsection}

% Macros
%
\def\lim{\mathop{\mathrm{lim}}\nolimits}
\def\colim{\mathop{\mathrm{colim}}\nolimits}
\def\Spec{\mathop{\mathrm{Spec}}}
\def\Hom{\mathop{\mathrm{Hom}}\nolimits}
\def\Ext{\mathop{\mathrm{Ext}}\nolimits}
\def\SheafHom{\mathop{\mathcal{H}\!\mathit{om}}\nolimits}
\def\SheafExt{\mathop{\mathcal{E}\!\mathit{xt}}\nolimits}
\def\Sch{\mathit{Sch}}
\def\Mor{\mathop{\mathrm{Mor}}\nolimits}
\def\Ob{\mathop{\mathrm{Ob}}\nolimits}
\def\Sh{\mathop{\mathit{Sh}}\nolimits}
\def\NL{\mathop{N\!L}\nolimits}
\def\CH{\mathop{\mathrm{CH}}\nolimits}
\def\proetale{{pro\text{-}\acute{e}tale}}
\def\etale{{\acute{e}tale}}
\def\QCoh{\mathit{QCoh}}
\def\Ker{\mathop{\mathrm{Ker}}}
\def\Im{\mathop{\mathrm{Im}}}
\def\Coker{\mathop{\mathrm{Coker}}}
\def\Coim{\mathop{\mathrm{Coim}}}

% Boxtimes
%
\DeclareMathSymbol{\boxtimes}{\mathbin}{AMSa}{"02}

%
% Macros for moduli stacks/spaces
%
\def\QCohstack{\mathcal{QC}\!\mathit{oh}}
\def\Cohstack{\mathcal{C}\!\mathit{oh}}
\def\Spacesstack{\mathcal{S}\!\mathit{paces}}
\def\Quotfunctor{\mathrm{Quot}}
\def\Hilbfunctor{\mathrm{Hilb}}
\def\Curvesstack{\mathcal{C}\!\mathit{urves}}
\def\Polarizedstack{\mathcal{P}\!\mathit{olarized}}
\def\Complexesstack{\mathcal{C}\!\mathit{omplexes}}
% \Pic is the operator that assigns to X its picard group, usage \Pic(X)
% \Picardstack_{X/B} denotes the Picard stack of X over B
% \Picardfunctor_{X/B} denotes the Picard functor of X over B
\def\Pic{\mathop{\mathrm{Pic}}\nolimits}
\def\Picardstack{\mathcal{P}\!\mathit{ic}}
\def\Picardfunctor{\mathrm{Pic}}
\def\Deformationcategory{\mathcal{D}\!\mathit{ef}}
\usepackage{cleveref}
\crefname{originallemma}{original lemma}{original lemmas}
\Crefname{originallemma}{Original lemma}{Original lemmas}
\makeatletter
\newlabel{lemma-exact-support-coherent}{{lemma-exact-support-coherent}{source}{}{sourcepointer-1}{}}
\newlabel{lemma-exact-support-coherent@cref}{{[originallemma][0][]lemma-exact-support-coherent}{[1][0][]source}{}{}{}}
\makeatother
\crefname{originalreference}{original reference}{original references}
\Crefname{originalreference}{Original reference}{Original references}
\makeatletter
\newlabel{more-algebra-lemma-dual-perfect-complex}{{more-algebra-lemma-dual-perfect-complex}{source}{}{sourcepointer-2}{}}
\newlabel{more-algebra-lemma-dual-perfect-complex@cref}{{[originalreference][0][]more-algebra-lemma-dual-perfect-complex}{[1][0][]source}{}{}{}}
\makeatother
\crefname{originalreference}{original reference}{original references}
\Crefname{originalreference}{Original reference}{Original references}
\makeatletter
\newlabel{more-algebra-lemma-perfect-ring-map}{{more-algebra-lemma-perfect-ring-map}{source}{}{sourcepointer-3}{}}
\newlabel{more-algebra-lemma-perfect-ring-map@cref}{{[originalreference][0][]more-algebra-lemma-perfect-ring-map}{[1][0][]source}{}{}{}}
\makeatother
\title{Boundedness and coherence properties preserved by the upper shriek of a finite type homomorphism}
\author{The Stacks Project Authors}
\date{stacks-project:dualizing.tex:local-snapshot-20261012; GNU FDL 1.2 or later}
\begin{document}
\maketitle
\noindent\textit{Editorial note.} All delivered mathematical text is the author's own, copied from the pinned Stacks Project chapter source sources/stacks-dualizing/pinned-original-source.tex; no statement or proof step was rewritten, reordered or removed.\par\medskip
\noindent\textit{Editorial note.} Preamble provenance: the Stacks Project keeps one single shared preamble for every chapter; this card takes it directly from the staged sources/stacks-dualizing/preamble.tex as three verbatim spans (lines 8-12, 14-18 and 154-217), so no derived file was written and nothing was added or redefined. The wrapper class and the base packages are supplied by the assembler.\par\medskip
\noindent\textit{Editorial note.} Bibliography provenance: neither the counted statement nor its proof contains a citation (this lane rescans the whole span before writing the card), so the card carries no bibliography and no entry was reconstructed.\par\medskip
\noindent\textit{Editorial note.} Selection provenance: this card was opened under the file-based control contract (CONTROL.md), with the mechanical pre-card gate run on its statement and proof interval before the card was written.\par\medskip
\noindent\textit{Allowed standard prerequisites (editorial).}\begin{itemize}
\item The preservation of bounded below coherent complexes by the derived internal hom along a finite ring map. Invoked inside the counted proof (source lines 4646-4648). Supplies the bounded below and coherent case for the internal hom factor of the shriek.
\item The description of the restriction of the internal hom as the tensor with a perfect complex in the perfect case. Invoked inside the counted proof (source lines 4652-4654). Identifies the internal hom with the tensor product with a perfect complex whose representatives are finite complexes of finite projectives.
\item The perfectness of a module over a polynomial ring in the perfect case of a finite type homomorphism. Invoked inside the counted proof (source lines 4649-4650). Provides the hypothesis that makes the tensor description available.
\end{itemize}\medskip
\ifdefined\hypertarget\hypertarget{sourcepointer-1}{}\fi
\noindent\textit{Source reference.} Original lemma lemma-exact-support-coherent: sources/stacks-dualizing/pinned-original-source.tex, source lines 2480--2480. Navigation pointer only: the counted text refers to the Stacks item lemma-exact-support-coherent, which is an existing labeled result of this same chapter and is not part of this card. The pointer carries the target's own label as its number and its definition line as its span; its content is not reproduced here.\par\medskip
\ifdefined\hypertarget\hypertarget{sourcepointer-2}{}\fi
\noindent\textit{Source reference.} Original reference more-algebra-lemma-dual-perfect-complex: sources/stacks-dualizing/pinned-original-source.tex, source lines 4653--4653. Navigation pointer only: the counted text refers to the Stacks item more-algebra-lemma-dual-perfect-complex, which belongs to another chapter of the project and is not part of this card. The pointer carries the target's own label as its number because this chapter's staged snapshot does not contain the other chapter's tag table; its content is not reproduced here.\par\medskip
\ifdefined\hypertarget\hypertarget{sourcepointer-3}{}\fi
\noindent\textit{Source reference.} Original reference more-algebra-lemma-perfect-ring-map: sources/stacks-dualizing/pinned-original-source.tex, source lines 4650--4650. Navigation pointer only: the counted text refers to the Stacks item more-algebra-lemma-perfect-ring-map, which belongs to another chapter of the project and is not part of this card. The pointer carries the target's own label as its number because this chapter's staged snapshot does not contain the other chapter's tag table; its content is not reproduced here.\par\medskip
\begin{lemma}
\label{lemma-shriek-boundedness}
Let $\varphi : R \to A$ be a finite type homomorphism of Noetherian rings.
\begin{enumerate}
\item $\varphi^!$ maps $D^+(R)$ into $D^+(A)$ and
$D^+_{\textit{Coh}}(R)$ into $D^+_{\textit{Coh}}(A)$.
\item if $\varphi$ is perfect, then $\varphi^!$ maps
$D^-(R)$ into $D^-(A)$,
$D^-_{\textit{Coh}}(R)$ into $D^-_{\textit{Coh}}(A)$, and
$D^b_{\textit{Coh}}(R)$ into $D^b_{\textit{Coh}}(A)$.
\end{enumerate}
\end{lemma}

\begin{proof}
Choose a factorization $R \to P \to A$ as in the definition of $\varphi^!$.
The functor $- \otimes_R^\mathbf{L} : D(R) \to D(P)$ preserves
the subcategories
$D^+, D^+_{\textit{Coh}}, D^-, D^-_{\textit{Coh}}, D^b_{\textit{Coh}}$.
The functor $R\Hom(A, -) : D(P) \to D(A)$
preserves $D^+$ and $D^+_{\textit{Coh}}$ by
Lemma \ref{lemma-exact-support-coherent}.
If $R \to A$ is perfect, then $A$ is perfect as a $P$-module, see
More on Algebra, Lemma \ref{more-algebra-lemma-perfect-ring-map}.
Recall that the restriction of $R\Hom(A, K)$ to $D(P)$ is
$R\Hom_P(A, K)$. By More on Algebra, Lemma
\ref{more-algebra-lemma-dual-perfect-complex}
we have $R\Hom_P(A, K) = E \otimes_P^\mathbf{L} K$ for
some perfect $E \in D(P)$. Since we can represent $E$ by
a finite complex of finite projective $P$-modules
it is clear that $R\Hom_P(A, K)$ is in
$D^-(P), D^-_{\textit{Coh}}(P), D^b_{\textit{Coh}}(P)$
as soon as $K$ is. Since the restriction functor
$D(A) \to D(P)$ reflects these subcategories, the
proof is complete.
\end{proof}

\end{document}
